% Integrar bajo 4.1.1. El detalle de matrices se entrega en el archivo de anexos.
\subsubsection{Modelo de amenazas y decisiones de seguridad}
\label{sec:4-seguridad-rt}

La configuración de seguridad \texttt{SEC-1.0}, fechada el 6 de octubre de 2026, desarrolla los controles de la Figura~\ref{fig:4-6} sobre los componentes y las fronteras efectivas de la solución. El Encargado de Seguridad mantiene su versión; el Arquitecto de Solución comprueba las interfaces y el Líder de Calidad verifica los criterios de aceptación. La configuración constituye una especificación contractual de construcción y operación. Sus registros de ejecución se producirán durante el contrato, conforme al criterio de acreditación de BT~1.5, sin atribuir a esta oferta una instalación, prueba de intrusión o auditoría ya realizada \parencite[secc.~1.5, p.~4; RT-11.01--06, p.~23]{pucv-transversal}.

Se utiliza STRIDE: suplantación, alteración, repudio, divulgación, denegación de servicio y elevación de privilegios. El análisis recorre cada componente, sus datos, sus entradas y salidas, y el cruce de fronteras entre Internet, nube, recinto, dispositivo, administración y tercero. Se considera tanto un atacante externo como una credencial o dispositivo comprometidos y un uso indebido autorizado. La matriz A4.1, incluida en \texttt{SYNAPTIX-Subdocumento4-Anexos}, identifica amenaza, control, responsable y comprobación por componente e integración; también relaciona cada control con ISO/IEC~27001:2022, anexo~A, e ISO/IEC~27002:2022. STRIDE organiza escenarios de abuso; Zero Trust determina cómo se comprueba cada solicitud y cómo se reduce su alcance \parencites{sec-stride,sec-iso27001,sec-iso27002,nist-zero-trust}.

Las decisiones se mantienen consistentes con las autoridades de 4.1 y SD5. Una alerta, un sensor o una réplica no autoriza una maniobra ni modifica por sí sola una deuda. Una baja conocida domina una autorización anterior; una ocupación o situación escolar indeterminada se hace visible y exige el procedimiento operacional correspondiente. Los archivos aportados por terceros no se ejecutan ni pasan directamente a evidencia aceptada. El aislamiento técnico evita perder simultáneamente el par local, y la restitución de servicio conserva los hechos durables y su conciliación.

La matriz se revisa antes de H3, antes de cada liberación y ante un cambio de región, componente, versión mayor, permiso, tratamiento de datos o integración. Un cambio conserva el registro anterior, motivo, activos afectados, alternativa evaluada, decisión, control modificado y evidencia exigida. Seguridad propone el tratamiento; Arquitectura y el dueño del proceso verifican sus efectos; Calidad bloquea una liberación cuyo riesgo crítico o alto siga abierto. La revisión mensual de controles no sustituye esta revisión por cambio.

\subsubsection{Clasificación de información y custodia criptográfica}
\label{sec:4-clasificacion-rt}

El Líder de Datos asigna clasificación a campos, documentos, eventos, exportaciones y copias derivadas; el responsable del proceso del CLIENTE aprueba finalidad y acceso. La Tabla~\ref{tab:4-clasificacion-rt} sintetiza cuatro niveles. La matriz completa A4.1 aplica estos niveles al tránsito, almacenamiento, consulta, exportación, conservación y eliminación.

\begin{table}[htbp]
\centering
\caption{Niveles de clasificación y restricciones diferenciales}\label{tab:4-clasificacion-rt}
\begin{tabular}{L{\dimexpr.16\linewidth-2\tabcolsep\relax}L{\dimexpr.27\linewidth-2\tabcolsep\relax}L{\dimexpr.29\linewidth-2\tabcolsep\relax}L{\dimexpr.28\linewidth-2\tabcolsep\relax}}
\hline
\EncabezadoTabla{Nivel} & \EncabezadoTabla{Información} & \EncabezadoTabla{Acceso} & \EncabezadoTabla{Control adicional} \\\hline
Pública & Información aprobada de servicios & Lectura publicada & Aprobación de publicación \\\hline
\rowcolor{SynSurface}Interna & Configuración no secreta y operación agregada & Personal según función & Exportación autorizada \\\hline
Confidencial & Contratos, deuda y documentación operacional & Rol y relación con el expediente & Consulta y exportación auditadas \\\hline
\rowcolor{SynSurface}Restringida & Salud, menores, personas y posición consentida de C07 & Finalidad, relación y MFA & Cifrado de campo y minimización local \\\hline
\end{tabular}
\FuenteTabla{elaboración propia; clasificación y cifrado conforme a \parencites[RT-11.03/09/10, p.~23]{pucv-transversal}[cap.~15, RT-11.10/RT-16.09, pp.~33--35]{caso07}.}
\end{table}

El nivel público permite divulgar únicamente contenido expresamente aprobado; no elimina los controles de integridad. Todos los niveles se cifran en reposo. La deuda exige auditoría de consulta aunque su clasificación no la convierta en un campo de salud. La categoría restringida comprende todos los datos personales de armadores, tripulantes y trabajadores contratistas y todos los datos personales de menores, además de salud y posición compartida con consentimiento. Las copias locales, índices, conjuntos analíticos y exportaciones heredan las restricciones de sus datos de origen.

En AWS se seleccionan claves simétricas administradas por el CLIENTE en KMS, separadas por ambiente y finalidad: datos transaccionales, documentos, auditoría, respaldos y claves de datos de campos. Se activa rotación automática de material cada 365 días. La rotación de KMS conserva material anterior para lectura; no recifra automáticamente el historial. Las claves de datos de campo se cambian cada 90 días y ante compromiso, con recifrado controlado de los datos activos y preservación de versiones necesarias para la retención. Una clave asimétrica de firma de procedencia se sustituye mediante una nueva clave anual y doble período de verificación; no se le atribuye la rotación automática de una clave simétrica \parencite{sec-kms-rotacion}.

AWS KMS protege las claves de datos AES-256-GCM descritas en 4.1 y SD5. El contexto de cifrado vincula ambiente, dominio, clasificación e identificador del objeto, sin incluir valores personales. Cada valor usa nonce único de 96 bits y etiqueta de autenticación de 128 bits; se bloquea reutilizar un nonce con la misma clave. RDS, S3, archivos temporales, estados de OpenTofu, registros y copias se cifran; los volúmenes locales utilizan LUKS2 y el parque Windows conserva BitLocker con Sophos Device Encryption. Android protege almacenamiento y la base operacional de la aplicación; el subconjunto restringido conserva además el cifrado de campo. Las claves locales permanecen en una bóveda cifrada vinculada al TPM del nodo, con copia de recuperación bajo doble custodia y acceso sólo del servicio autorizado. Esta custodia no prolonga permisos offline ni modifica el RLH.

El CLIENTE aprueba custodios y política. Seguridad administra políticas de claves; SRE ejecuta recuperación y rotación autorizadas; la identidad de aplicación obtiene sólo el uso criptográfico de su finalidad. Quien administra una clave no recibe por esa función permiso de lectura de datos. Deshabilitar o destruir material requiere doble autorización, identificación de todas las copias afectadas y comprobación de restauración; se impide eliminar una clave necesaria para información retenida. CloudTrail registra cambios y uso de claves; las exportaciones registran solicitante, finalidad, alcance y manifiesto sin reproducir contenido sensible en los registros técnicos.

La matriz A4.1 aplica ISO/IEC~27017:2026 a responsabilidades compartidas, segregación de ambientes, administración, redes virtuales, supervisión y salida del servicio. Para ISO/IEC~27018:2025 desarrolla finalidad del tratamiento, instrucciones del CLIENTE, subencargados, ubicaciones de procesamiento, minimización, transparencia, acceso y eliminación. AWS administra infraestructura y servicios administrados; Synaptix configura permisos, datos, software, claves y operación; el CLIENTE decide finalidad y usuarios de negocio. Los servicios adicionales no heredan automáticamente la ubicación de RDS. La habilitación de una integración exige documentar su ubicación efectiva y tratamiento contractual, incluido el control de acceso del prestador \parencites{sec-iso27017-actual,sec-iso27018-actual}.

\subsubsection{Exposición, contratos de API y archivos entrantes}
\label{sec:4-exposicion-rt}

La publicación conserva CloudFront, AWS WAF, Shield Standard y API Gateway REST API Regional. El inventario A4.1 distingue nombres públicos de la oferta, endpoints administrados, rutas de terceros y servicios privados; una dirección de origen alcanzable técnicamente queda inventariada aunque su política rechace el acceso directo. Se deshabilita \texttt{execute-api}, se rechaza un origen sin el encabezado secreto de CloudFront y se restringen rutas, métodos y destinos. El ALB, RDS, EMQX, BMC y puertos de administración carecen de ingreso público. Los únicos puertos públicos de aplicación son TCP~443 y, para redirección a HTTPS, TCP~80 en el ingreso web. Las VPN del recinto admiten únicamente UDP~500/4500 entre pares autorizados, sin consola administrativa por Internet.

Se exige TLS~1.3 en clientes y terminaciones bajo control de la solución y se prohíben TLS~1.0/1.1 y SSL. CloudFront utiliza certificado propio de ACM y política \texttt{TLSv1.3\_2025}; REST API Regional utiliza \texttt{SecurityPolicy\_TLS13\_1\_3\_2025\_09}. El listener HTTPS del ALB usa \texttt{ELBSecurityPolicy-TLS13-1-3-2021-06}; los destinos propios sólo aceptan TLS~1.3. Se configura HTTPS exclusivo al origen y se verifica la negociación de cada tramo, incluidos Cognito, SDK, servicios AWS, sincronización y cada tercero. Se admiten suites TLS AES-GCM y ChaCha20-Poly1305 modernas. Las capacidades de un tramo no se infieren de las del ingreso: el inventario registra protocolo y excepción técnica de la contraparte, y una integración que no sostenga el requisito no se habilita mediante una rebaja silenciosa \parencites{sec-cloudfront-tls,sec-api-tls,sec-alb-tls}.

Las respuestas web incluyen \texttt{Strict-Transport-Security: max-age=63072000; includeSubDomains; preload}. El plazo equivale a $2\times365\times24\times3600$ segundos. Se configura redirección a HTTPS y todos los subdominios antes de solicitar su inclusión en la lista de precarga; la directiva no equivale a una inclusión ya obtenida. ACM renueva certificados administrados y se supervisan estados y vencimientos. Certificados importados y certificados del recinto utilizan renovación automática con doble versión temporal, prueba de cadena y reemplazo; se alerta a 60, 30, 15 y 7 días del vencimiento. SRE responde a un fallo de renovación el mismo día, sin esperar la última alerta \parencite{sec-hsts}.

El Lambda authorizer verifica identidad, emisor, audiencia, dispositivo, vigencia y versión de sesión; el dominio vuelve a comprobar autorización de recurso y revocación. La caché del authorizer permanece deshabilitada para acceso operacional y privilegiado. El modelo de solicitud y la validación de aplicación rechazan esquema, formatos, campos adicionales, tipo, longitud y reglas no verificadas por API Gateway. Cada método limita JSON a 64~KiB, rechaza sobredimensión en WAF y utiliza el contrato OpenAPI~3.1 con la proyección compatible ya descrita. El inventario de cuotas A4.1 fija tasas, ráfagas, unidades y contadores de efecto; una API key nunca sustituye autenticación. Los planes de uso de AWS limitan abuso por esfuerzo máximo, mientras una cuota de negocio se comprueba mediante un contador transaccional antes de producir su efecto \parencites{p3-api-cuotas,p3-api-validacion,p3-waf-cuerpos}.

Se selecciona GuardDuty Malware Protection for S3 para la cuarentena regional, con SSE-KMS, rol limitado al prefijo de entrada y resultados por EventBridge. La aplicación verifica tamaño, firma binaria, tipo admitido, huella, propietario, manifestación de consentimiento y pertenencia del objeto al expediente. Sólo \texttt{NO\_THREATS\_FOUND}, obtenido para la misma versión y huella, permite su promoción; \texttt{THREATS\_FOUND}, \texttt{UNSUPPORTED}, \texttt{ACCESS\_DENIED}, \texttt{FAILED} y ausencia de resultado mantienen el objeto aislado y generan incidencia. No se aceptan ejecutables, macros ni archivos protegidos por contraseña. La falta de hallazgo no acredita autenticidad documental. Se aceptan inicialmente JPEG/PNG y PDF sin contenido activo, hasta 20~MiB por objeto, con inspección adicional de estructura y un manifiesto de cantidad/tamaño. Este límite es de diseño y no una cuota de catálogo \parencites{sec-gdu-s3,sec-gdu-s3-estados,sec-gdu-s3-cifrado}.

El registro durable de un hecho operacional no espera el examen remoto de su adjunto. La cola conserva el original cifrado y la asociación; los hechos válidos y el lote completo se sincronizan dentro del plazo del caso. El adjunto no se publica, procesa mediante Textract ni consume como evidencia aceptada antes del examen. La prueba conjunta debe incluir demora y fallo de examen dentro de la ventana de treinta minutos; separar cuarentena de publicación no autoriza diferir el envío obligatorio del archivo.

Para el portal se incorpora AWS WAF Bot Control con reglas dirigidas a login, recuperación, alta y consultas repetitivas. La secuencia de respuesta es observación, limitación, reto silencioso y reto accesible cuando el riesgo persista; se bloquea una explotación confirmada. Se conservan alternativas mediante autenticación reforzada y asistencia verificada, con teclado, lector de pantalla, tiempo suficiente y recuperación de operación. Los retos de navegador no se aplican a callbacks firmados, sincronización o telemetría: esas rutas usan identidad técnica, firma, límites y deduplicación. La prueba de accesibilidad comprueba también falso positivo y salida del reto. No se permite que una regla por IP compartida corte toda la operación de la marina \parencites{sec-waf-bot,sec-waf-retos}.

\subsubsection{Vulnerabilidades y cobertura de detección y respuesta}
\label{sec:4-vulnerabilidades-rt}

El registro \texttt{VUL-1.0} vincula cada activo de T-11 con responsable, modelo, sistema, versión, firmware, dependencia, exposición, herramienta de descubrimiento y fuente de avisos. La cobertura inicial comprende 21 estaciones Windows, dos nodos Ubuntu, 16 terminales administrados y todas las tareas, imágenes y funciones cloud, además de redes, BMC, UPS, climatización, acceso, video e instrumentación. Las reservas enroladas también reciben controles. Un equipo sin agente admisible conserva supervisión de red, inventario de firmware y revisión automatizada de avisos del fabricante; no se le atribuye un agente Windows o Linux incompatible.

Inspector se configura para escaneo mejorado continuo de todos los repositorios ECR y para funciones Lambda soportadas; la duración de reexamen se fija para conservar cobertura de imágenes en uso y cada despliegue comprueba elegibilidad. Trivy~0.74.0 examina dependencias, imágenes y archivos de infraestructura en cada cambio y recompone hallazgos con base de vulnerabilidades actualizada diariamente. En Ubuntu se inventarían diariamente paquetes de distribución, Pro y contenido de contenedores; avisos de Canonical y fabricantes se cotejan al menos cada cuatro horas. Intune entrega inventario y cumplimiento de actualizaciones de las 21 estaciones; sus paquetes y aplicaciones se cotejan con los avisos de Microsoft y su fabricante, sin atribuirle por ello un escáner EDR. Sophos Mobile aporta versiones y postura; las aplicaciones propias se revisan mediante su SBOM y Android mediante el boletín de seguridad y soporte del modelo. Los activos sin herramienta activa registrada no cuentan como cubiertos \parencites{sec-inspector,sec-trivy,p3-intune-licencia}.

Cada hallazgo conserva primera publicación y primera detección; el vencimiento se calcula desde el instante anterior aplicable y nunca se reinicia al duplicarlo, cambiar su responsable o reabrirlo. Los máximos son siete días corridos para crítico, quince para alto y treinta para medio. El escaneo y la ingestión continua detectan cambios; la conciliación diaria comprueba fuentes ausentes y activos no observados. Seguridad triagea y asigna; SRE o el dueño de software corrige; Calidad valida por reexamen y caso de explotación o regresión. Si no existe parche, se retira, sustituye o deshabilita efectivamente la función vulnerable antes del máximo y se comprueba la eliminación de exposición. Una aceptación interna de riesgo no amplía un plazo de las bases \parencite[RT-11.04, p.~23]{pucv-transversal}.

La protección mantiene las familias seleccionadas: Sophos Endpoint con XDR y Device Encryption para Windows; Sophos Protection for Linux~2026.3 con XDR--Server para ambos nodos; Sophos Mobile MDM+MTD e Intercept~X for Mobile~9.8.4146 para la flota; y GuardDuty Runtime Monitoring con sidecar para cada tarea Fargate soportada. En Windows, la línea inicial Core Agent~2026.2.2.1.0 procede de las notas oficiales de septiembre de 2026. Los servicios SaaS se identifican por edición, tenant, configuración \texttt{SEC-1.0} y exportación de política; las versiones reales de sus agentes se incorporan al registro de liberación y se actualizan bajo soporte \parencites{sec-sophos-windows,p3-sophos-linux,sec-sophos-android,p3-guardduty}.

Se habilitan antimanipulación, detección en ejecución, telemetría, cuarentena y acceso de respuesta limitado a Seguridad/SRE. GuardDuty dirige su hallazgo a EventBridge, SOC y la guía de contención: preservar identificadores y registros, retirar tráfico de la tarea afectada, detenerla cuando exista réplica sana, bloquear el digest y volver a una imagen aceptada. En Linux se preserva evidencia y se verifica el nodo sano antes de aislar el otro. En Windows se aísla el dispositivo y se coordina continuidad del puesto; en Android se bloquea acceso de negocio mediante la Autoridad local y se aplica la acción móvil admitida al recuperar canal. MTD móvil no se presenta como EDR de servidor y el análisis de imágenes no sustituye la detección en ejecución.

Cada despliegue compara activo esperado con agente, política, licencia, último contacto y capacidad de respuesta. El recinto emite latidos de cobertura cada minuto y alerta tras cinco minutos sin protección local; una pérdida WAN se clasifica como falta de telemetría remota, no como ausencia automática del agente. La consola central y los terminales tienen reglas de contacto propias por misión y turno. Agente detenido, política desactivada, licencia ausente o imagen sin análisis impiden habilitar un activo nuevo; la operación existente se contiene conforme al riesgo y se mantiene su procedimiento de continuidad. Las exclusiones exigen archivo/proceso exacto, motivo, vencimiento y revisión de Seguridad; una exclusión global de PostgreSQL, Podman o todo el sistema queda prohibida.

\subsubsection{SIEM, centro de seguridad y respuesta a incidentes}
\label{sec:4-siem-incidentes-rt}

Se selecciona Sophos XDR Powered by Secureworks con el complemento Sophos Next-Gen SIEM, edición SaaS, para correlacionar eventos propios y de terceros sin desplegar otra plataforma en el recinto. La documentación de la edición exige licenciar el complemento sobre los usuarios y servidores XDR/MDR y ofrece 13 meses de retención; se usa el tenant proyectado Brasil/São Paulo de Sophos. La geografía efectiva del servicio de correlación y su tratamiento contractual se individualizan junto con la suscripción antes de habilitar envío, sin inferirlos del tenant de protección de endpoints. Se preservan las regiones AWS primaria y secundaria y la exclusión de OpenSearch \parencites{sec-sophos-siem,sec-sophos-tenant}.

Se normalizan eventos de negocio, identidad, gateway, WAF, CloudTrail, GuardDuty, Security Hub, red, Linux, Windows, móviles y cambios de configuración. El ingreso cloud utiliza HTTP Ingest autenticado desde el proceso de Auditoría y seguridad, después de persistir la evidencia en S3; las fuentes locales se almacenan cifradas y se reenvían con identificador y secuencia. Parsers propios convierten eventos al esquema soportado y conservan el identificador de origen. El enriquecimiento correlaciona sujeto seudónimo, dispositivo, expediente seudónimo, turno, política y estado de negocio antes de publicar una señal; no se presupone que el lenguaje de consultas del producto pueda resolver cualquier unión entre tablas \parencites{sec-sophos-ingesta,sec-sophos-reglas}.

Los casos específicos A4.1 detectan consulta o exportación indebida de salud y menores, retiro escolar incompatible, cambio de permiso después de una baja, alteración de despacho de combustible, duplicación de cobro, modificación de evidencia ambiental, cambio de posición sin consentimiento, manipulación de un aviso de marejada, repetición o secuencia inválida de eventos y falta de auditoría. Se configuran reglas y agrupación por incidente, no sólo alarmas de CPU. Los registros técnicos omiten datos sensibles y credenciales; la evidencia detallada de una actuación permanece cifrada y restringida en el dominio, con una referencia opaca en el SIEM.

S3 en una cuenta separada de Seguridad/Registros conserva los lotes originales con versionado y Object Lock en modo compliance, manifiesto de huellas y cadena de custodia. Los eventos de seguridad se mantienen doce meses en almacenamiento en línea y veinticuatro meses adicionales en archivo recuperable, para un total de 36 meses; se usa calendario mensual y fecha de evento, no una aproximación de treinta días por mes. El archivado conserva las claves y la capacidad de recuperación. Los trece meses del SIEM complementan la búsqueda y no sustituyen la copia inalterable ni autorizan acortar la conservación. Se comprueban cada mes recuperación de un lote, huellas, permisos y continuidad de la cadena \parencite[RT-11.14, p.~24]{pucv-transversal}.

El SOC se presta desde el centro remoto de Centinela Austral Servicios TI SpA en Santiago, bajo supervisión del Encargado de Seguridad de Synaptix, coherente con SD1. Se reserva un puesto de analista activo 24\,$\times$\,7, supervisor de escalamiento y especialistas de Synaptix de disponibilidad. Para dimensionar la reserva, una jornada de planificación de 40~h/semana y 20\,\% de indisponibilidad para descanso, formación y ausencias proporciona $40(1-0,20)=32$~h/persona/semana; cubrir $24\times7=168$~h exige $\lceil168/32\rceil=6$ personas en el pool de turnos. Es una capacidad mínima a contratar para un puesto continuo y no seis personas simultáneas, ni una atribución de personal ya asignado. El supervisor es una función adicional de disponibilidad que no consume las horas del puesto activo. La dotación, suplencias y agenda quedan separadas de los ocho integrantes propios de SD1 y de la mesa de usuarios.

El plan \texttt{IR-1.0} de A4.1 aplica NIST SP~800-61 revisión~3. Desde la primera detección se abre un registro con reloj UTC, severidad, activo, alcance, evidencia, medidas y responsables. Crítico comprende riesgo para personas por manipulación de escuela, marejada o maniobra, compromiso privilegiado activo, ransomware o extracción de datos restringidos; alto comprende intrusión localizada o pérdida relevante de control sin ese efecto confirmado; medio y bajo conservan dueño y seguimiento. El analista clasifica en un máximo de quince minutos y escala inmediatamente el crítico a Seguridad y SRE. La Contraparte Técnica y el representante de gerencia del CLIENTE reciben llamada y notificación escrita dentro de dos horas de la detección de un crítico, aun si el alcance sigue bajo investigación \parencites{sec-nist-incidentes}[RT-11.18, p.~24]{pucv-transversal}.

Toda brecha de seguridad o de datos personales, incluso no crítica, se comunica al CLIENTE dentro de veinticuatro horas de su detección mediante informe preliminar: instantes, hechos conocidos, datos y personas potencialmente afectados, alcance confirmado y no confirmado, contención, continuidad, evidencia conservada y próxima actualización. No se exige confirmar todos los hechos para iniciar ese reloj. El análisis de causa raíz se entrega dentro de cinco días hábiles siguientes, con causas, cronología, controles fallidos y acciones con dueño/plazo; las acciones se siguen hasta cierre verificado. La obligación de aviso contractual no sustituye las comunicaciones regulatorias que correspondan al responsable del tratamiento \parencite[RT-11.19, p.~24; RT-14.06, p.~27]{pucv-transversal}.

Los procedimientos conservan copia forense o exportación íntegra admisible, huellas, autorización de adquisición, custodio y accesos. No se borra ni reinicia un equipo antes de la preservación necesaria, salvo medida inmediata justificada para seguridad de personas. La recuperación usa configuración conocida y verifica datos, permisos, sincronización y controles antes de volver a servicio. El informe evita reproducir salud, credenciales o datos de menores en avisos generales. Un relevo de turno entrega incidentes abiertos, relojes, acciones y destinatarios, con aceptación del siguiente analista.

\subsubsection{Construcción segura, promoción y cadena de suministro}
\label{sec:4-devsecops-rt}

La política \texttt{CI-1.0} convierte los controles de entrega en puertas automáticas. Git del CLIENTE conserva requerimientos, incidencias, código, contratos, infraestructura y migraciones; la rama principal y las de liberación prohíben escritura directa, borrado y force-push. Cada cambio usa solicitud de integración con revisión de una persona distinta del autor, aprobación adicional del dueño de seguridad para políticas/secretos y revalidación cuando cambia el commit. El rol de integración automática sólo fusiona un commit aprobado con todos los controles exitosos; una persona administradora no dispone de una excepción permanente para omitirlos.

El manifiesto de liberación vincula ID de requisito/RF/RNF/RT, incidencia, solicitud, commit y árbol resueltos, versión de contrato/esquema, resultados de prueba, SBOM, digest, firma, procedencia, ambiente, despliegue y reversión. Una prueba conserva versión, entorno, datos, instante, resultado y evidencia. El pipeline bloquea una incidencia sin requerimiento, una liberación sin prueba válida para el mismo digest o una promoción que recompila el artefacto. QA, Preproducción y Producción reciben exactamente el mismo artefacto; configuración y secretos se inyectan por ambiente.

Se especifican CodeBuild y CodePipeline administrados en AWS para construcción y promoción. Las herramientas iniciales son Semgrep~1.179.0 para análisis estático de Java, Kotlin y TypeScript/JavaScript; Trivy~0.74.0 para composición, infraestructura, imágenes y generación de SBOM CycloneDX; Gitleaks~8.30.1 para secretos e historial; ZAP~2.17.0 para análisis dinámico autenticado en QA/Preproducción; y Cosign~3.1.3 para firma y verificación. Se comprueban publicación oficial, firma/huella de distribución, licencia y dependencias de la propia herramienta; se conservan digest y revisiones en la BOM. La versión anterior Cosign~3.0.5 no se adopta porque el fabricante documenta fallos de verificación corregidos en 3.1.3. Tampoco se utilizan tags móviles de acciones ni la distribución Trivy~0.69.4 comprometida. Las herramientas se actualizan antes de su pérdida de soporte y conforme a los máximos de vulnerabilidades \parencites{sec-semgrep,sec-trivy,sec-gitleaks,sec-zap,sec-cosign,sec-cosign-aviso}.

La Tabla~\ref{tab:4-ci-puertas-rt} muestra los mínimos bloqueantes. Los análisis diarios reevalúan las dependencias aceptadas; no se considera seguro un commit sólo porque una base de vulnerabilidades antigua no produjo hallazgos.

\begin{table}[htbp]
\centering
\caption{Puertas automáticas de la configuración CI-1.0}\label{tab:4-ci-puertas-rt}
\begin{tabular}{L{\dimexpr.27\linewidth-2\tabcolsep\relax}L{\dimexpr.39\linewidth-2\tabcolsep\relax}L{\dimexpr.34\linewidth-2\tabcolsep\relax}}
\hline
\EncabezadoTabla{Control} & \EncabezadoTabla{Condición de bloqueo} & \EncabezadoTabla{Evidencia} \\\hline
Construcción y pruebas & Fallo o cobertura de negocio menor a 70\,\% & Reporte ligado al commit \\\hline
\rowcolor{SynSurface}SAST, SCA y contenedores & Hallazgo crítico/alto sin cierre & Reportes y base de análisis \\\hline
Secretos & Una credencial confirmada o prueba incompleta & Hallazgo redaccionado y rotación \\\hline
\rowcolor{SynSurface}DAST y autorización & Crítico/alto o acceso indebido & Informe autenticado y roles \\\hline
Artefacto y procedencia & Firma, digest o builder inválidos & SBOM y atestación verificada \\\hline
\end{tabular}
\FuenteTabla{elaboración propia; \parencite[RT-04.05/11, pp.~10--11; RT-11.22--27, p.~24; RT-20.01, p.~34]{pucv-transversal}.}
\end{table}

El umbral de cobertura se calcula como líneas de lógica de negocio ejecutadas divididas por líneas ejecutables de esa lógica. Se identifica previamente el conjunto de código de los nueve dominios; excluir un caso de uso para alcanzar el porcentaje queda prohibido. El pipeline agrega los reportes sin duplicar líneas y falla ante falta de reporte. El análisis estático utiliza reglas versionadas para autorización, inyección, exposición de datos, cifrado, concurrencia e idempotencia; DAST incluye usuarios con roles distintos, sesión revocada y entradas adversas. Críticos y altos deben cerrarse antes de producción conforme a RT-20.01, además de respetar los plazos de remediación.

El gestor Secrets Manager rota automáticamente las credenciales de aplicaciones PostgreSQL cada treinta días con alternancia de usuarios y prueba de lectura/escritura antes de promoción; las de integración y el encabezado de origen se sustituyen cada treinta días mediante función de rotación con doble versión breve y revocación de la anterior. Identidades AWS de cargas usan roles y credenciales temporales, sin claves estáticas. La bóveda local ejecuta el mismo ciclo sobre secretos locales; una falla de enlace no convierte un secreto vencido en indefinido. La rotación modifica tanto el almacén como el sistema que acepta la credencial y alerta si no puede completar comprobación o revocación. Nunca se inyecta un secreto como texto dentro de una imagen, repositorio o archivo de configuración persistente \parencite{sec-secrets-rotacion}.

El diseño de procedencia \texttt{PROV-1.0} exige SLSA Build~L3 conforme a la especificación~1.2: construcción hospedada y efímera por ejecución, ausencia de caché mutable compartida, aislamiento de ejecuciones concurrentes y generación de atestación por un plano de control separado. Ese generador obtiene commit resuelto, configuración y parámetros de la API/control de ejecución, enumera parámetros externos y verifica digest del artefacto recibido. Los pasos de construcción no tienen permiso de firma KMS ni de modificación del generador; una identidad exclusiva firma la procedencia y Cosign verifica confianza en builder, commit, parámetros, firma y digest antes de promover. ECR conserva tags inmutables y la política de despliegue admite únicamente digest autorizado. La firma sola no establece L3 ni se atribuye ese nivel automáticamente a CodeBuild: la calificación de la plataforma debe demostrar que el código del build no puede falsificar la atestación, acceder a la clave o influir otra ejecución \parencites{sec-slsa,sec-codebuild}.

Los cinco ambientes estarán habilitados y operativos como condición de H3. Preproducción conserva topología, versiones y configuración productivas; sus diferencias de capacidad se registran y no eliminan controles o redundancia. La promoción se elige azul-verde en nube con enrutamiento inicial del 10\,\% durante cinco minutos, seguido de 100\,\% y observación de quince minutos; el pipeline vuelve automáticamente al digest anterior si falla una comprobación crítica, se incumple un SLO del flujo afectado o falla integridad de datos. Antes de cada paso a producción se ejecuta y registra esa secuencia completa en Preproducción, incluyendo reversión. En el par local se actualiza primero el pasivo, se verifica y conmuta bajo fencing, manteniendo compatibilidad; no se actualizan ambos nodos simultáneamente. Los períodos propuestos son parámetros de aceptación iniciales y se validan con la carga del caso, sin reducir las puertas contractuales.

Las migraciones de esquema se versionan y se automatizan mediante scripts PostgreSQL de incorporación, transición y retirada. Cada versión define migración directa y retorno: se agrega estructura compatible con ambas versiones, se migra por lotes idempotentes y se verifica conciliación; la estructura anterior y los datos originales se conservan hasta cerrar la ventana de retorno. Antes de un retiro irreversible se conserva una copia verificable y se prueba restauración y reaplicación de hechos de la ventana. El pipeline impide una transformación destructiva sin camino de retorno y prueba de convivencia de las dos versiones de aplicación. La reversión de aplicación no se confunde con recuperar datos destruidos.

La aprobación de una dependencia nueva requiere origen oficial, versión y digest, licencia compatible, mantenimiento activo, soporte durante su uso, compatibilidad del stack y ausencia de vulnerabilidades conocidas sin corregir. El solicitante aporta SBOM y pruebas; Seguridad verifica vulnerabilidades y procedencia; Arquitectura revisa compatibilidad; Calidad verifica regresión y la revisión competente aprueba licencia. Se deniega una dependencia de origen no comprobable o con exposición conocida sin eliminar. Un registro aprobado conserva fecha, revisores por función y nueva revisión ante cambio de versión, licencia o fin de soporte. La evaluación OWASP SAMM~2.2.0 se efectúa antes de H3 y se repite anualmente durante Operación, con evidencia por práctica y un plan de mejora medible, sin atribuir un nivel de madurez inicial aún no evaluado \parencite{sec-samm}.

La prueba de intrusión se contratará a un tercero independiente de Synaptix, de su equipo de construcción y del SOC operador. Se realizará anualmente y antes de cada paso a producción, incluyendo cada liberación que se promueva; alcance, credenciales de prueba y reglas de intervención cubren portal, API, nube, nodo local, móvil, integraciones y autorización. El CLIENTE recibirá el informe íntegro, sin sustitución por resumen ejecutivo, y el plan de remediación con plazos y reexamen. Críticos/altos bloquean el paso. Cada año de Operación se ejecuta además un ejercicio de incidente con participación del CLIENTE, escenario de menor/marejada y pérdida de enlace, tiempos de aviso, roles, observaciones y acciones. No se ejecuta explotación intrusiva durante una ventana operacional protegida \parencite[RT-11.20/21, p.~24]{pucv-transversal}.

\subsubsection{Acceso técnico y evidencias de operación}
\label{sec:4-acceso-alertas-rt}

Las personas desarrolladoras carecen de acceso interactivo directo a Producción, bases o shell de las tareas. El acceso excepcional usa incidencia, finalidad, activo, autorización de Seguridad y responsable del proceso, MFA, elevación por un máximo de sesenta minutos, caducidad automática y grabación íntegra de sesión. Se ejecuta desde el canal de administración autorizado con comandos permitidos y exportación protegida. AWS Systems Manager Session Manager registra sesiones de shell admitidas; se prohíben túneles de puerto y SSH a través de Session Manager porque AWS no registra su contenido. ECS Exec o una modalidad sin grabación verificable no se habilitan como atajo. Para el nodo local se emplea la sesión remota grabada por el servicio de acceso técnico y un agente de captura local; su identidad temporal y autorización también se comprueban sin Internet. Si no puede grabarse, se deniega el acceso excepcional y se utiliza el procedimiento de recuperación autorizado \parencites{sec-ssm-grabacion}[RT-11.27, p.~24]{pucv-transversal}.

El alertamiento comparte identificador de transacción y evidencia con OpenTelemetry y CloudWatch, pero agrega síntomas de negocio. Un retiro escolar no difundido, una amarra sin estado confirmado, un plan vencido sin acuse de escalamiento o una campaña de marejada con destinatarios pendientes genera una incidencia aunque CPU y red estén disponibles. Las reglas A4.1 declaran condición, ventana, gravedad, fuente, agrupación, responsable y guía. Se agrupan duplicados por servicio, expediente seudónimo, causa y ventana de cinco minutos; se conserva cada evento original. Las dependencias suprimen avisos subordinados de infraestructura, pero no eliminan alertas de seguridad de personas, falta de auditoría o vencimiento contractual. Un mantenimiento autorizado silencia sólo los activos/horario declarados.

El SOC atiende seguridad y el NOC atiende disponibilidad, con recepción 24\,$\times$\,7 y escalamiento a Seguridad/SRE. Una incidencia sin aceptación en cinco minutos escala al supervisor; al cumplirse quince minutos escala a los líderes Synaptix y se comprueba contacto por canal alternativo. Las guías de A4.1 recogen pérdida WAN, nodo, quorum/fencing, datos, colas, sincronización, reloj/permisos, certificados/claves, malware, proveedor externo, alertas de negocio y energía/clima/incendio. Cada una define disparador, diagnóstico seguro, acción permitida, parada, reversión, verificación, comunicación y evidencias. Las tareas repetibles se automatizan en ese orden después de probarse, sin automatizar decisiones marítimas ni la eliminación de evidencia.

Un incidente crítico de cualquier origen produce análisis de causa raíz dentro de cinco días hábiles, con acciones seguidas hasta cierre. El CLIENTE conserva acceso propio a tableros y exportaciones, con la restricción de información de cada perfil. Las métricas técnicas y de negocio se retienen doce meses en línea y veinticuatro en archivo de agregados; las trazas y registros técnicos, treinta días en línea y once meses en archivo, salvo que integren un incidente o una obligación mayor. Los eventos de seguridad mantienen su regla separada de doce más veinticuatro meses. El inventario de observabilidad identifica volumen, ingestión, consultas, recuperación y duplicación; su valorización se desarrolla en la Oferta Económica y no en este subdocumento, conforme al Art.~50.2. La memoria económica y el dimensionamiento completo de consumo deben corresponder al volumen efectivo antes de dar por cerrado RT-14.08 \parencite[RT-14.01--08, p.~27]{pucv-transversal}.

La aceptación verifica configuraciones exportadas, negativa de acceso, cifrado y rotación, cobertura perdida, alteración de procedencia, contención de un nodo sin detener el par, examen de archivos fallido, tiempos de aviso y restauración de evidencia. El archivo A4.1 contiene los criterios y productos de comprobación por control. La firma del acta se sustenta en resultados obtenidos para las versiones desplegadas; la matriz T-12 distingue esa evidencia de ejecución de la formulación verificable contenida en esta oferta.
